\section{What the certificates do and what remains open}

The practical output is an explicitly bounded comparison for an actual,
infinite-volume AQ state in a small nonzero-coupling patterned subfamily. The
reference curve is calculable and the interval pays the difference from that
reference; this does not turn the reference into the exact interacting answer.
A common enclosing interval for different allowed subsequential states is not
a proof that those states coincide. Deterministic analytic proxies are not
independent observations, and their errors cannot be divided by the square
root of the number of nodes.

The direct proof avoids a quantitative finite-volume error modulus for its own
analytic enclosure: local reset and dynamics estimates hold uniformly before
the already admitted AQ limit. It does not certify an arbitrary finite-box
solver. Such a solver still needs state/boundary comparison, complete
representation-cutoff and omitted-channel bounds, actual ground projection and
energy errors, evolution error, centering error, arithmetic error and the
specified physical clock. Observed stability of digits does not supply these
missing hypotheses.

The unchanged fundamental limits are substantial: no uniform Wilson magnetic
Hamiltonian is proved by the zero-selected calculation; no all-state uniqueness,
AO/AQ identification, quantitative boundary independence, resolved particle
pole, dispersion law, susceptibility or matched continuum theory follows.
An inverse-energy form is mathematically defined on the vacuum-orthogonal
sector; identifying it with a thermodynamic response derivative would require
a separately specified deformation and differentiability theorem. The Clay
problem requires the four-dimensional continuum construction and its physical
properties, not only a fixed-spacing estimate \cite{clay}.

The updated
\repo{research/round31/advisor/roadmap.json}{prospective roadmap}
orders the next questions by what this work does not establish:
\begin{enumerate}
\item Return to $|\tau|=10^{-8}$ with a sharper local state or observable
estimate. A creation-expansion approach must pay its complete normalization,
support and remainder costs; a formal leading term is insufficient.
\item Resolve an interaction-induced effect by an interval that excludes its
matched free value, with the observable and separation target fixed in advance.
\item Match selected and omitted coefficients to a uniform Hamiltonian,
charging selected interactions or proving a genuinely non-Haar reference
comparison. Even a uniform strong-coupling result is not a continuum limit.
\item Compare declared boundary and state constructions quantitatively in the
required local topology and at a common physical clock.
\item Establish a regulator/coupling trajectory, scale control and the
continuum reconstruction hypotheses, or isolate a specific necessary estimate
that fails along the chosen route.
\end{enumerate}
Every item is unexecuted. No fraction of the final continuum proof is inferred
from the number of loops, equations, pages or named contributions. The existing
historical-panel, Newton, Tesla and paired-research skills were retained; future
investigations require fresh prospective contracts and source snapshots.
